\documentclass[a4paper]{article}

\usepackage[margin=1in]{geometry}
\usepackage{amsmath,amsfonts,mathtools,amsthm,amssymb}
\usepackage[l2tabu,orthodox]{nag}
\usepackage{microtype} % fixes spacing or whatever
\usepackage{enumitem}
\usepackage{tikz-cd}
\usepackage{xcolor}
\usepackage{physics}
\usepackage{mathrsfs}
\usepackage{cancel}
% \usepackage{libertine,libertinust1math}
\usepackage{eucal}
\usepackage[toc,page]{appendix}
\usepackage{pgfplots}
\pgfplotsset{compat=1.18}

\newtheorem{proposition}{Proposition}
\newtheorem{theorem}{Theorem}
\newtheorem{lemma}{Lemma}
\newtheorem{corollary}{Corollary}

\theoremstyle{definition}
\newtheorem{remark}{Remark}
\newtheorem{definition}{Definition}
\newtheorem{example}{Example}

% \usepackage[sc,osf]{mathpazo}   % With old-style figures and real smallcaps.
% \linespread{1.025}              % Palatino leads a little more leading
% % Euler for math and numbers
% \usepackage[euler-digits,small]{eulervm}
% \AtBeginDocument{\renewcommand{\hbar}{\hslash}}

\usepackage{etoolbox}
\usepackage{dynkin-diagrams}
\def\row#1/#2!{\dynkin{#1}{#2}\\}
\newcommand{\tble}[1]{
   \renewcommand*\do[1]{\row##1!}
   \[
      \begin{array}{ll}\docsvlist{#1}\end{array}
   \]
}

\allowdisplaybreaks

\renewcommand{\epsilon}{\varepsilon}
% \renewcommand{\phi}{\varphi}

\DeclareMathAlphabet{\mathpzc}{OT1}{pzc}{m}{it}

\newcommand{\goth}[1]{\mathfrak{#1}}
\newcommand{\red}[1]{#1_{\text{red}}}
\newcommand{\ad}[1]{#1_{\text{Ad}}}
\newcommand{\reg}[1]{#1_{\text{reg}}}
\newcommand{\sh}[1]{#1^{\text{sh}}}
\newcommand{\cpdv}[2]{\cfrac{\partial #1}{\partial #2}}

\newcommand{\fA}{\mathcal{A}}
\newcommand{\fB}{\mathcal{B}}
\newcommand{\fC}{\mathcal{C}}
\newcommand{\fD}{\mathcal{D}}
\newcommand{\fE}{\mathcal{E}}
\newcommand{\fF}{\mathcal{F}}
\newcommand{\fG}{\mathcal{G}}
\newcommand{\fH}{\mathcal{H}}
\newcommand{\fI}{\mathcal{I}}
\newcommand{\fJ}{\mathcal{J}}
\newcommand{\fK}{\mathcal{K}}
\newcommand{\fL}{\mathcal{L}}
\newcommand{\fM}{\mathcal{M}}
\newcommand{\fN}{\mathcal{N}}
\newcommand{\fO}{\mathcal{O}}
\newcommand{\fP}{\mathcal{P}}
\newcommand{\fQ}{\mathcal{Q}}
\newcommand{\fR}{\mathcal{R}}
\newcommand{\fS}{\mathcal{S}}
\newcommand{\fT}{\mathcal{T}}
\newcommand{\fX}{\mathcal{X}}
\newcommand{\fY}{\mathcal{Y}}
\newcommand{\fZ}{\mathcal{Z}}

\newcommand{\be}{\mathbf{e}}
\newcommand{\bx}{\mathbf{x}}
\newcommand{\bone}{\mathbf{1}}
\newcommand{\cx}{\bar{x}}
\newcommand{\cy}{\bar{y}}
\newcommand{\vx}{\vec{x}}
\newcommand{\vsigma}{\vec{\sigma}}
\newcommand{\tzeta}{\tilde{\zeta}}

\newcommand{\gm}{\goth{m}}
\newcommand{\gp}{\goth{p}}
\newcommand{\gF}{\goth{F}}
\newcommand{\gU}{\goth{U}}
\newcommand{\gV}{\goth{V}}

\newcommand{\A}{\mathbf{A}}
\newcommand{\C}{\mathbf{C}}
\newcommand{\F}{\mathbf{F}}
\newcommand{\G}{\mathbf{G}}
\newcommand{\bH}{\mathbf{H}}
\newcommand{\N}{\mathbf{N}}
\newcommand{\PP}{\mathbf{P}}
\newcommand{\Q}{\mathbf{Q}}
\newcommand{\R}{\mathbf{R}}
\newcommand{\Z}{\mathbf{Z}}

\newcommand{\dlog}{\dd\log}

\newcommand\srestr[2]{{\left.\kern-\nulldelimiterspace #1\vphantom{\small|} \right|_{#2}}}
\newcommand\restr[2]{{\left.\kern-\nulldelimiterspace #1 \vphantom{\big|} \right|_{#2}}}
\newcommand\gsec[2]{\Gamma({#1}, {#2})}
\newcommand\gsecx[1]{\Gamma(X, {#1})}

\DeclareMathOperator{\chr}{char}
\DeclareMathOperator{\rProj}{\mathbf{Proj}}
\DeclareMathOperator{\rSpec}{\mathbf{Spec}}
\DeclareMathOperator{\rHom}{\mathbf{Hom}}
\DeclareMathOperator{\rExt}{\mathbf{Ext}}
\DeclareMathOperator{\id}{id}
\DeclareMathOperator{\Frac}{Frac}
\DeclareMathOperator{\rk}{rank}
\DeclareMathOperator{\deter}{det}
\DeclareMathOperator{\pic}{Pic}
\DeclareMathOperator{\cacl}{CaCl}
\DeclareMathOperator{\trd}{tr.d.}
\DeclareMathOperator{\cl}{Cl}
\DeclareMathOperator{\depth}{depth}
\DeclareMathOperator{\codim}{codim}
\DeclareMathOperator{\Div}{Div}
\DeclareMathOperator{\coker}{coker}
\DeclareMathOperator{\len}{length}
\DeclareMathOperator{\height}{height}
\DeclareMathOperator{\supp}{Supp}
\DeclareMathOperator{\proj}{Proj}
\DeclareMathOperator{\im}{im}
\DeclareMathOperator{\Hom}{Hom}
\DeclareMathOperator{\Der}{Der}
\DeclareMathOperator{\spec}{Spec}
\DeclareMathOperator{\Aut}{Aut}
\DeclareMathOperator{\ch}{char}
\DeclareMathOperator{\tor}{Tor}
\DeclareMathOperator{\Ann}{Ann}
\DeclareMathOperator{\Syl}{Syl}
\DeclareMathOperator{\Sym}{Sym}
\DeclareMathOperator{\GL}{GL}
\DeclareMathOperator{\SL}{SL}
\DeclareMathOperator{\Stab}{Stab}
\DeclareMathOperator{\Perm}{Perm}
\DeclareMathOperator{\Orb}{Orb}
\DeclareMathOperator{\Gal}{Gal}
\DeclareMathOperator{\Supp}{Supp}
\DeclareMathOperator{\Ext}{Ext}
\DeclareMathOperator{\Tor}{Tor}
\DeclareMathOperator{\PGL}{PGL}
\DeclareMathOperator{\hd}{hd}
\DeclareMathOperator{\pd}{pd}
\DeclareMathOperator{\SO}{SO}
\DeclareMathOperator{\SU}{SU}
\DeclareMathOperator{\Sp}{Sp}
\DeclareMathOperator{\Oth}{O}
\DeclareMathOperator{\Uni}{U}
\DeclareMathOperator{\evf}{ev}
\DeclareMathOperator{\cH}{H}
\DeclareMathOperator{\dR}{R}
\DeclareMathOperator{\End}{End}
\DeclareMathOperator{\Hasse}{Hasse}
\DeclareMathOperator{\Num}{Num}

\author{James Lee}

% Solarized
% \pagecolor[RGB]{0,20,26}
% \color[RGB]{191,191,191}

% Sephia
% \pagecolor[RGB]{249,239,220}

% Grey
% \pagecolor[RGB]{200, 200, 200}

% Black
% \pagecolor[RGB]{0,0,0}
% \color[RGB]{255,255,255}


\title{Chapter 3, Section 10}

\begin{document}

\maketitle

\begin{enumerate} [label=\textbf{\arabic*.}, leftmargin=0em]

\item Over a nonperfect field, smooth and regular are not equivalent.
For example, let $k_0$ be a field of characteristic $p > 0$, let $k = k_0(t)$, and let $X \subseteq \A_k^2$ be the curve defined by $y^2 = x^p - t$.
Show that every local ring of $X$ is a regular local ring, but $X$ is not smooth over $k$.

\begin{proof}
  To show every local ring of $X$ is a regular local ring, it suffices to show the coordinate ring of $X$ defined by $A = k[x, y] / (y^2 - x^p - t)$ is integrally closed.
  If $p \neq 2$, the result follows from (II, Ex. 6.4), so assume $p = 2$.
  An immediate consequence is
  \begin{equation*}
    y^2 - x^2 - t = (x + y)^2 - t,
  \end{equation*}
  which implies the isomorphism
  \begin{equation*}
    A = \frac{k[x, y]}{(x + y)^2 - t} \cong k(\sqrt{t})[z].
  \end{equation*}
  But $k(\sqrt{t})$ is a field, and polynomial rings over fields are integrally closed, so we conclude that $A$ is integrally closed.

  Let $\bar{k}$ be the algebraic closure of $k$.
  By (10.2), $X$ is smooth over $k$ if and only if $X \times_k \bar{k}$ is regular, i.e., $A \otimes_k \bar{k}$ defines a nonsingular $\bar{k}$-variety; but $t^{1/p} \in \bar{k}$, so we can factorize
  \begin{equation*}
    y^2 - x^p - t = y^2 - (x + t^{1/p})^p.
  \end{equation*}
  Thus, $(-t^{1/p}, 0)$ is a singular point of $X \times_k \bar{k}$ (A.M., Ex. 11.1), so $X$ is not smooth over $k$.
\end{proof}

\item Let $f : X \to Y$ be a proper, flat morphism of varieties over $k$. Suppose for some point $y \in Y$ that the fiber $X_y$ is smooth over $k(y)$. Then show that there is an open neighborhood $V$ of $y$ in $Y$ such that $f : f^{-1}(V) \to V$ is smooth.

\begin{proof}
  Flat morphisms of finite type between Noetherian schemes are open (Ex. 9.1), proper morphisms are closed, $X$ and $Y$ are irreducible, so $f$ must be surjective, and the generic point $\xi \in X$ is mapped to the generic point of $Y$.
  Some immediate consequence are the following: $f$ trivially satisfies (2) in the definition of smoothness of relative dimension $n = \dim{X} - \dim{Y}$, so for every open set $V \subseteq X$, $f : f^{-1}(V) \to V$ satisfies (2).
  Thus, we want to find an open neighborhood $V$ of $y$ such that $\Omega_{X / Y}$ is free at every point in $f^{-1}(V)$ (10.0.2).
  For every $x \in X_y$, we have the inequalities
  \begin{equation*}
    \dim{X} - \dim{Y} \leq \dim_{k(\xi)} \Omega_{X / Y} \otimes k(\xi) \leq \dim_{k(x)} \Omega_{X / Y} \otimes k(x) = \dim{X_y},
  \end{equation*}
  by (II, Ex. 5.8a) and (II, 8.6A), where the last equality is by smoothness of $X_y$ over $k(y)$ and the fact that relative differentials commute with base change.
  Since $n = \dim{X_y}$ (9.6), $\Omega_{X / Y}$ is free of rank $n$ at every $x \in X_y$ (II, 8.9).

  Now, let $U = \{ x \in X \mid \Omega_{X / Y} \text{ is free at $x$}\}$, and $Z = X - U$.
  Being free is an open property for finitely generated modules, so $U$ is an open subset of $X$.
  Then $f$ is an open and closed mapping, so $f(U)$ and $X - f(Z)$ are open sets in $Y$.
  Also, $f^{-1}(y) = X_y$ is a subset of $U$, which means $y$ is not in $f(Z)$.
  Therefore, $f(U) \cap (X - f(Z))$ is a nonempty open set containing $y$, so we can find an open neighborhood $V$ of $y$ such that $V \subseteq f(U) \cap (X - f(Z))$.
  Since $f^{-1}(V) \subseteq U$, $\Omega_{X / Y}$ is free at every point in $f^{-1}(V)$, which is what we wanted to show.
\end{proof}

\item A morphism $f : X \to Y$ of schemes of finite type over $k$ is \textit{étale} if it is smooth of relative dimension $0$. It is \textit{unramified} if for every $x \in X$, letting $y = f(x)$, we have $\goth{m}_y \cdot \fO_{x, X} = \goth{m}_x$, and $k(x)$ is a separable algebraic extension of $k(y)$. Show that the following conditions are equivalent:
\begin{itemize}
  \item[(i)] $f$ is étale;
  \item[(ii)] $f$ is flat, and $\Omega_{X/Y} = 0$;
  \item[(iii)] $f$ is flat and unramified.
\end{itemize}

\begin{proof}
  $(i) \iff (ii)$ The forward direction is trivial.
  On the other hand, $\Omega_{X/Y} = 0$ trivially satisfies condition (3) in the definition of smoothness, and we want to show condition (2) holds as well.
  According to (9.6), condition (2) is equivalent to saying that every irreducible component of every fibre $X_y$ of $f$ has dimension $0$, so let $X'$ be an irreducible component of $X_y$ with its reduced induced structure.
  Then $X'$ is an integral scheme of finite type over $k(y)$, and $\Omega_{X'/k(y)}$ is a coherent sheaf, so by (II, Ex. 3.20b) and (II, 8.6A),  we reduce to showing $\Omega_{X'/k(y)} \otimes k(x) = 0$ for every $x \in X'$.
  Then
  \begin{equation*}
    \Omega_{X_y / k(y)} \otimes k(x) = \Omega_{X / Y} \otimes k(x) = 0
  \end{equation*}
  since relative differentials commute with base extension and $\Omega_{X / Y} = 0$.
  Hence, the result follows from the surjective map of (II, 8.12)
  \[ \begin{tikzcd}
    \Omega_{X_y/k(y)} \otimes k(x) \arrow[r] & \Omega_{X'/k(y)} \otimes k(x) \arrow[r] & 0.
  \end{tikzcd} \]

  $(i) \iff (iii)$ Relative differentials commute with base extension, so by the equivalence of (i) and (ii), we reduce to the case when $Y = \spec{k}$.
  Thus, we want to show $X$ is smooth of relative dimension 0 over $k$ if and only if $\goth{m}_x = 0$ and $k(x)$ is a separable algebraic extension of $k$ for every $x \in X$.
  Consider the exact sequence of (II, 8.4A)
  \[ \begin{tikzcd}
    \goth{m}_x/\goth{m}_x^2 \arrow[r] & \Omega_{X/k} \otimes k(x) \arrow[r] & \Omega_{k(x)/k} \arrow[r] & 0.
  \end{tikzcd} \]
  In the forward direction, suppose $\Omega_{X / k} = 0$.
  Then $\Omega_{k(x) / k} = 0$ by the exact sequence, so $k(x)$ is a separable algebraic extension of $k$ (II, 8.6A)
  Also, $X$ is smooth of relative dimension 0 over $k$ if and only if $\bar{X} = X \times \bar{k}$ is equidimensional of dimension zero and regular, where $\bar{k}$ an algebraic closure of $k$ (10.2).
  Regularity and dimension of $X \times \bar{k}$ implies $\goth{m}_x \otimes_k \bar{k} = 0$, hence $\goth{m}_x = 0$.
  On the otherhand, if $\goth{m}_x = 0$ and $k(x)$ is a separable algebraic extension of $k$.
  Then $\Omega_{k(x) / k} = 0$ (II, 8.6A), and the term on the left hand of the exact sequence is zero.
  Hence, $\Omega_{X / k} = 0$.
\end{proof}

\item[\textbf{5.}] If $x$ is a point of a scheme $X$, we define an \textit{étale neighborhood} of $x$ to be an étale morphism $f : U \to X$, together with a point $x' \in U$, such that $f(x') = x$. As an example of the use of étale neighborhoods, prove the following: if $\fF$ is a coherent sheaf on $X$, and if every point of $X$ has an étale neighborhood $f : U \to X$ for which $f^*\fF$ is a free $\fO_U$-module, then $\fF$ is locally free on $X$.

\begin{proof}
  Consider the following lemma:
  \begin{lemma}
    If $B$ is a faithfully flat $A$-algebra, then for any $A$-module $M$, we have $M$ is $A$-flat if and only if $M \otimes_A B$ is $B$-flat.
  \end{lemma}
  Let $\goth{a}$ be an ideal of $A$.
  Then
  \begin{equation*}
    \tor_1^A(M, A / \goth{a}) \otimes_A B \cong \tor_1^B(M \otimes_A B, B / \goth{a}B) = 0,
  \end{equation*}
  so by faithfull flatness, we conclude that $\tor_1^A(M, A / \goth{a}) = 0$ (A.M., Ex. 3.16).

  Now, if $f^*\fF$ is a free $\fO_U$-module, then $(f^*\fF)_{x'}$ is a free $\fO_{x', U}$-module with $(f^*\fF)_{x'} \cong \fF_x \otimes_{\fO_{x, X}} \fO_{x', U}$, and $\fO_{x', U}$ is a faithfully flat $\fO_{x, X}$-algebra, so the result follows from Lemma 1 and the fact that for finitely generated modules over Noetherian local rings, being free is equivalent to being flat.
\end{proof}

\item[\textbf{6.}] Let $Y$ be the plane nodal cubic curve $y^2 = x^2(x + 1)$. Show that $Y$ has a finite étale covering $X$ of degree 2, where $X$ is a union of two irreducible components, each one isomorphic to the normalization of $Y$.

\begin{proof}
  For simplicity, assume $k$ is algebraically closed and has characteristic $\neq 2$.
  Consider two copies of the twisted cubic curve $X_1$ and $X_2$ in $\A^3$ with affine coordinates $x, y, z$ given by the parametric equations
  \begin{equation*}
    X_1 : \begin{cases}
      x = t^2 - 1 \\
      y = t^3 - t \\
      z = t
    \end{cases}, \quad X_2 : \begin{cases*}
      x = t^2 - 1 \\
      y = t^3 - t \\
      z = -t
    \end{cases*}.
  \end{equation*}
  Each are isomorphic to the normalization of $Y$, and they intersect at the following three points
  \begin{equation*}
    P_1 = (0, 0, 1), \quad P_2 = (0, 0, -1), \quad Q = (-1, 0, 0).
  \end{equation*}
  Let $X$ be the closed subscheme of $X$ with two irreducible components $X_1$ and $X_2$.
  There is an obvious map $X \to Y$ given by the projection map onto the $xy$-plane, and each irreducible component of $X$ maps onto $Y$.
  Finally, let $\pi : \tilde{X} \to X$ be obtained from $X$ by blowing up the point $Q$.
  Then $\tilde{X} \to Y$ is our desired étale covering of degree 2.
\end{proof}

\end{enumerate}
\end{document}

% \item[\textbf{7.}] \textit{(Serre). A linear system with moving singularities.} Let $k$ be an algebraically closed field of characteristic $2$. Let $P_1, \dots, P_7 \in \PP_k^2$ be the seven points of the projective plane over the prime field $\F_2 \subseteq k$. Let $\goth{d}$ be the linear system of all cubic curves in $\PP^2$ passing through $P_1, \dots, P_7$.
% \begin{itemize}
%   \item[(a)] $\goth{d}$ is a linear system of dimension $2$ with base points $P_1, \dots, P_7$, which determines an inseparable morphism of degree $2$ from $\PP^2 - \{P_i\}$ to $\PP^2$.
%   \item[(b)] Every curve $C \in \goth{d}$ is singular. More precisely, either $C$ consists of $3$ lines all passing through one of the $P_i$, or $C$ is an irreducible cuspidal cubic with cusp $P \neq \text{any $P_i$}$. Furthermore, the correspondence $C \mapsto \text{the singular point of $C$}$ is a one-to-one correspondence between $\goth{d}$ and $\PP^2$. Thus, the singular points of elements of $\goth{d}$ move all over.
% \end{itemize}
% 
% \begin{proof} $ $ \vspace{0pt}
%   \begin{itemize} [leftmargin=0cm]
%     \item[(a)] 
% 
%     \item[(b)]
%   \end{itemize}
% \end{proof}


% \item Show that a morphism $f : X \to Y$ of schemes of finite type over $k$ is étale if and only if the following condition is satisfied: for each $x \in X$, let $y = f(x)$. Let $\hat{\fO}_x$ and $\hat{\fO}_y$ be the completions of the local rings at $x$ and $y$. Choose fields of representatives (II, 8.25A) $k(x) \subseteq \hat{\fO}_x$ and $k(y) \subseteq \hat{\fO}_y$ so that $k(y) \subseteq k(x)$ via the natural map $\hat{\fO}_y \to \hat{\fO}_x$. Then our condition is that for every $x \in X$, $k(x)$ is a separable algebraic extension of $k(y)$, and the natural map
% \begin{equation*}
%   \hat{\fO}_y \otimes_{k(y)} k(x) \to \hat{\fO}_x
% \end{equation*}
% is an isomorphism.



% \item \textit{A linear system with moving singularities contained in the base locus (any characteristic).} In affine $3$-space with coordinates $x, y, z$, let $C$ be the conic $(x - 1)^2 + y^2 = 1$ in the $xy$-plane, and let $P$ be the point $(0, 0, t)$ on the $z$-axis. Let $Y_t$ be the closure in $\PP^3$ of the cone over $C$ with vertex $P$. Show that as $t$ varies, the surfaces $\{Y_t \}$ form a linear system of dimension $1$, with a moving singularity at $P$. The base locus of this linear system is the conic $C$ plus the $z$-axis.

% \item Let $f : X \to Y$ be a morphism of varieties over $k$. Assume that $Y$ is regular, $X$ is Cohen-Macaulay, and that every fiber of $f$ has dimension equal to $\dim{X} - \dim{Y}$. Then $f$ is flat.